%% ============================================================
%%  Predicción Bayesiana del FIFA World Cup 2026
%%  Cesar Prieto — UBITS
%%  Compilar con: pdflatex → pdflatex (dos pasadas para referencias)
%% ============================================================
\documentclass[12pt,a4paper]{article}

%% --- Codificación y fuentes ---
\usepackage[utf8]{inputenc}
\usepackage[T1]{fontenc}
\usepackage[spanish,es-nodecimaldot]{babel}

%% --- Matemáticas ---
\usepackage{amsmath}
\usepackage{amssymb}
\usepackage{amsthm}

%% --- Tipografía y espaciado ---
\usepackage{microtype}
\usepackage{setspace}
\onehalfspacing
\usepackage{parskip}

%% --- Geometría ---
\usepackage[margin=2.5cm]{geometry}

%% --- Tablas ---
\usepackage{booktabs}
\usepackage{array}
\usepackage{multirow}
\usepackage{caption}
\usepackage[table]{xcolor}
\definecolor{lightgray}{gray}{0.93}

%% --- Figuras y flotantes ---
\usepackage{graphicx}
\usepackage{float}

%% --- Referencias ---
\usepackage{natbib}

%% --- Hipervínculos ---
\usepackage{hyperref}
\hypersetup{
  colorlinks = true,
  linkcolor  = blue!70!black,
  citecolor  = green!50!black,
  urlcolor   = blue!70!black,
  pdftitle   = {Predicción Bayesiana de Resultados del FIFA World Cup 2026},
  pdfauthor  = {Cesar Prieto}
}

%% --- Comandos matemáticos personalizados ---
\newcommand{\E}{\mathbb{E}}
\newcommand{\Var}{\operatorname{Var}}
\newcommand{\iid}{\overset{\text{\tiny iid}}{\sim}}

%% --- Entornos ---
\newtheorem{proposition}{Proposición}
\newtheorem{remark}{Observación}

%% ============================================================
\begin{document}

%% ---- Portada -----------------------------------------------
\title{%
  \textbf{Predicción Bayesiana de Resultados del FIFA World Cup 2026}\\[6pt]
  \large Modelo Jerárquico de Poisson con Corrección de Dixon-Coles
}
\author{%
  Cesar Prieto\\[4pt]
  \normalsize UBITS\\
  \normalsize \href{mailto:cesarprietosar@gmail.com}{cesarprietosar@gmail.com}
}
\date{Junio de 2026}

\maketitle
\thispagestyle{empty}

%% ---- Abstract ----------------------------------------------
\begin{abstract}
\noindent
Se presenta un sistema de predicción probabilística para el FIFA World Cup 2026
basado en un modelo jerárquico bayesiano de Poisson con corrección de dependencia
de \citet{dixon1997modelling}.
El conjunto de datos comprende más de 47{,}000 encuentros internacionales desde 1872
(fuente: \textit{martj42/international\_results}), del que se seleccionan
1{,}935 partidos disputados desde 2010 entre equipos clasificados al torneo,
ponderados exponencialmente por antigüedad (vida media $= 365$ días) y por
importancia competitiva.
El predictor log-lineal incorpora efectos aleatorios jerárquicos de ataque $\alpha_k$
y defensa $\delta_k$ por equipo, un coeficiente de ventaja local $\eta$ y un
coeficiente de diferencial de Elo $\beta$.
La corrección de Dixon-Coles —parametrizada por $\rho$— ajusta las probabilidades
de resultados exactos bajos (0-0, 1-0, 0-1, 1-1), corrigiendo la sobreestimación
inherente al supuesto de independencia entre los goles de ambos equipos.
La inferencia se realiza con Stan mediante 4 cadenas MCMC (8{,}000 iteraciones
efectivas); todos los parámetros convergen con $\hat{R} \leq 1{,}002$.
El coeficiente de Elo arroja $\hat{\beta} = 0{,}510$
(IC\,95\%: $[0{,}316;\, 0{,}696]$), confirmando el poder predictivo del historial
de rendimiento histórico.
El valor estimado $\hat{\rho} = -0{,}071$ implica un incremento relativo de
aproximadamente 10--15\% en la probabilidad de empates 0-0 y 1-1.
La simulación de 100{,}000 escenarios de fase de grupos y 10{,}000 de rondas
eliminatorias sitúa a España como favorita con una probabilidad de título del
13{,}5\%, seguida de Argentina (9{,}2\%) y Francia (8{,}1\%).

\bigskip
\noindent\textbf{Palabras clave:} modelos de conteo bayesianos,
modelo jerárquico de Poisson, corrección de Dixon-Coles,
ratings Elo, predicción deportiva, fútbol internacional.
\end{abstract}

\newpage
\tableofcontents
\newpage

%% ==========================================================
\section{Introducción}
\label{sec:intro}
%% ==========================================================

La predicción de resultados deportivos es un problema de larga data en
estadística aplicada, con implicaciones prácticas en mercados de apuestas,
planificación deportiva y análisis táctico.
El fútbol internacional presenta desafíos particulares: los encuentros son poco
frecuentes por par de equipos, los marcadores son típicamente bajos y la
correlación entre los goles de ambos rivales viola el supuesto de independencia
que subyace a los modelos de Poisson univariados estándar.

El presente trabajo describe un sistema bayesiano completo para pronosticar el
FIFA World Cup 2026, el primero de la historia celebrado en tres países
(Estados Unidos, Canadá y México) con 48 equipos participantes.
La ampliación del torneo incrementa la complejidad combinatoria de la fase de
grupos —doce grupos de cuatro equipos con clasificación de los dos primeros y
ocho mejores terceros— y justifica una aproximación probabilística que cuantifique
la incertidumbre en cada etapa.

Las contribuciones centrales del trabajo son tres:
\begin{enumerate}
  \item Un modelo jerárquico de Poisson con efectos aleatorios de ataque y defensa
        que permite el intercambio de información entre equipos (regularización parcial),
        estimado en el marco bayesiano completo con propagación de incertidumbre.
  \item La corrección de dependencia de \citet{dixon1997modelling} integrada
        directamente en la verosimilitud del modelo Stan, con el parámetro $\rho$
        estimado del posterior junto con todos los demás parámetros.
  \item Ratings Elo calculados desde cero sobre la totalidad del historial disponible
        e incorporados como covariable en el predictor, resolviendo el sesgo de
        sobreestimación de equipos con muchas goleadas contra rivales débiles.
\end{enumerate}

El resto del documento se organiza como sigue.
La Sección~\ref{sec:related} revisa la literatura relevante.
La Sección~\ref{sec:data} describe las fuentes de datos y la construcción del
conjunto de entrenamiento.
La Sección~\ref{sec:elo} detalla el cálculo de los ratings Elo.
La Sección~\ref{sec:model} especifica el modelo estadístico completo.
La Sección~\ref{sec:inference} presenta los resultados de la inferencia bayesiana.
La Sección~\ref{sec:model-comparison} compara el modelo propuesto con alternativas.
La Sección~\ref{sec:simulation} describe la mecánica de simulación del torneo.
La Sección~\ref{sec:results} resume las probabilidades obtenidas.
La Sección~\ref{sec:discussion} discute las implicaciones y limitaciones.
La Sección~\ref{sec:conclusion} concluye.

%% ==========================================================
\section{Trabajo Relacionado}
\label{sec:related}
%% ==========================================================

La modelización estadística de resultados de fútbol arranca formalmente con
\citet{maher1982modelling}, quien propone representar los goles marcados por
cada equipo como variables de Poisson independientes con tasas que dependen de
la fuerza atacante y la debilidad defensiva.
Aunque simple, el modelo de Maher captura los aspectos fundamentales de la
distribución marginal de goles y establece la parametrización ataque-defensa
que ha dominado la literatura posterior.

\citet{dixon1997modelling} identifican el defecto central del enfoque de
independencia: la distribución conjunta del marcador exhibe una correlación
negativa en resultados bajos, con sobreabundancia de empates a cero y a uno.
Proponen el factor de corrección $\tau(\cdot)$ que ajusta únicamente las celdas
$(0,0)$, $(1,0)$, $(0,1)$ y $(1,1)$ mediante un único parámetro $\rho$,
manteniendo la tractabilidad analítica.
Este trabajo es considerado el punto de inflexión entre los modelos puramente
descriptivos y los sistemas de predicción operativos.

\citet{rue2000prediction} extienden el enfoque con una parametrización dinámica
de las fuerzas de ataque y defensa, actualizando los parámetros en el tiempo
mediante un filtro de Kalman aproximado, anticipando la necesidad de modelar
la no-estacionariedad del rendimiento deportivo.

\citet{karlis2003analysis} proponen el modelo de Poisson bivariado, en el que
ambas puntuaciones comparten un componente latente que induce correlación positiva.
Aunque más general, este modelo no resulta superior al de \citet{dixon1997modelling}
en calibración de resultados de baja puntuación según comparaciones publicadas
en ligas europeas.

En el ámbito bayesiano, \citet{baio2010bayesian} aplican un modelo jerárquico
a la liga italiana Serie A y \citet{groll2018prediction} predicen copas del mundo
mediante regresión de Poisson regularizada a nivel de equipo.
Ninguno de estos trabajos incorpora simultáneamente la corrección de Dixon-Coles
y ratings Elo como covariable en un único marco de inferencia bayesiana completa,
brecha que el presente trabajo cierra.

%% ==========================================================
\section{Fuentes de Datos y Construcción del Dataset}
\label{sec:data}
%% ==========================================================

\subsection{Historial internacional de resultados}

El conjunto de datos primario proviene del repositorio público
\texttt{martj42/international\_results} en GitHub \citep{martj42},
que recopila sistemáticamente resultados oficiales de selecciones nacionales
desde 1872. En su versión utilizada, el repositorio contiene
\textbf{47{,}000+ partidos} con las siguientes variables relevantes:

\begin{center}
\begin{tabular}{ll}
\toprule
\textbf{Variable}   & \textbf{Descripción} \\
\midrule
\texttt{date}       & Fecha del encuentro \\
\texttt{home\_team} & Equipo local \\
\texttt{away\_team} & Equipo visitante \\
\texttt{home\_score}& Goles del equipo local \\
\texttt{away\_score}& Goles del equipo visitante \\
\texttt{tournament} & Nombre del torneo \\
\texttt{neutral}    & Indicador de sede neutral \\
\bottomrule
\end{tabular}
\end{center}

Para la estimación del modelo se aplican tres filtros: (1)~partidos desde
el 1 de enero de 2010, cubriendo cuatro ciclos de Copa del Mundo; (2)~al menos
uno de los equipos clasificado al WC 2026; (3)~marcador completo disponible.
Estos criterios generan una muestra inicial de \textbf{7{,}366 partidos únicos}.
Para el entrenamiento del modelo se impone un cuarto filtro: \textit{ambos}
equipos deben ser clasificados al WC 2026 (criterio \textit{WC vs.\ WC}),
garantizando que los efectos aleatorios estén respaldados por observaciones
directas entre los 48 participantes.
El dataset de entrenamiento final comprende \textbf{1{,}935 partidos}
(3{,}870 filas en formato largo, una por equipo por partido).

\subsection{Equipos participantes y fixtures}

La lista de 48 selecciones clasificadas y el fixture completo se extrajeron por
\textit{data mining} de ESPN, capturando los identificadores numéricos de equipo
y los slugs textuales directamente desde los atributos \texttt{href} de los
enlaces de la página de clasificaciones. Las discrepancias de nomenclatura entre
fuentes (e.g., \textit{Czech Republic} en el historial vs.\ \textit{Czechia} en
ESPN) se resuelven mediante una tabla de correspondencia manual validada.

\subsection{Ponderación temporal y por torneo}

Un partido jugado hace cinco años contiene información relevante pero menos
precisa sobre el nivel actual de un equipo que uno disputado recientemente.
Se asigna a cada partido un peso mediante decaimiento exponencial:
\begin{equation}
  w_t = \exp\!\left(-\frac{\log 2}{\tau_{1/2}} \cdot \Delta t\right),
  \qquad \tau_{1/2} = 365 \text{ días},
  \label{eq:temporal-weight}
\end{equation}
donde $\Delta t$ es el número de días transcurridos desde el partido.
La importancia competitiva se refleja mediante un multiplicador:
\begin{equation}
  m_{\text{torneo}} = \begin{cases}
    2{,}0 & \text{fase final de Copa del Mundo,} \\
    1{,}5 & \text{eliminatorias mundialistas,} \\
    1{,}0 & \text{torneos competitivos continentales,} \\
    0{,}5 & \text{amistosos.}
  \end{cases}
  \label{eq:tournament-mult}
\end{equation}
El peso compuesto de cada observación es $p_i = w_{t,i} \times m_{\text{torneo},i}$.

%% ==========================================================
\section{Ratings Elo}
\label{sec:elo}
%% ==========================================================

\subsection{Motivación e incorporación al modelo}

Los efectos aleatorios $\alpha_k$ y $\delta_k$ capturan heterogeneidad latente
en el período de estimación (2010--actualidad, WC vs.\ WC).
Sin embargo, cuando dos equipos han disputado pocos partidos directos, estos
efectos pueden estar pobremente identificados.
Los ratings Elo proveen una señal histórica acumulada —derivada de la totalidad
del historial desde 1872, incluyendo partidos fuera del dataset de entrenamiento—
que complementa y estabiliza la estimación jerárquica.
Adicionalmente, el Elo resuelve el sesgo de sobreestimación de equipos con
altas tasas de gol contra rivales débiles en eliminatorias: la covariable
$\text{elo\_diff}$ explica estructuralmente la diferencia de nivel, liberando
los efectos aleatorios para capturar solo desviaciones idiosincrásicas.

En el predictor lineal, el diferencial de Elo normalizado se define como:
\begin{equation}
  \text{elo\_diff}_i = \frac{\text{Elo}_{\text{equipo}} - \text{Elo}_{\text{rival}}}{400}.
  \label{eq:elo-diff}
\end{equation}
La normalización por 400 es convencional y hace que la escala de
$\text{elo\_diff}$ sea aproximadamente $[-3, 3]$, compatible con los priors del modelo.

\subsection{Cálculo y actualización}

Los ratings Elo se calculan procesando en orden cronológico los 47{,}000+
partidos del repositorio completo, con todos los equipos iniciando en
$R = 1500$. La fórmula de actualización es estándar:
\begin{equation}
  R'_A = R_A + K \cdot G \cdot (S_A - E_A),
  \label{eq:elo-update}
\end{equation}
donde $S_A \in \{1, 0{,}5, 0\}$ es el resultado observado,
\begin{equation}
  E_A = \frac{1}{1 + 10^{(R_B - R_A - h)/400}}
  \label{eq:elo-expected}
\end{equation}
es la probabilidad esperada con ventaja de local $h = 100$ (o $h = 0$ en
sede neutral), $G$ es el multiplicador por diferencia de goles:
\begin{equation}
  G = \begin{cases}
    1{,}0 & |GD| \leq 1, \\
    1{,}5 & |GD| = 2, \\
    (11 + |GD|) / 8 & |GD| \geq 3,
  \end{cases}
  \label{eq:elo-g}
\end{equation}
y $K$ es el factor de importancia por tipo de torneo:
\begin{equation}
  K = \begin{cases}
    60 & \text{Copa del Mundo (fase final),} \\
    50 & \text{campeonatos continentales,} \\
    40 & \text{eliminatorias mundialistas,} \\
    35 & \text{torneos competitivos,} \\
    20 & \text{amistosos.}
  \end{cases}
  \label{eq:elo-k}
\end{equation}

%% ==========================================================
\section{Modelo Estadístico}
\label{sec:model}
%% ==========================================================

\subsection{Predictor log-lineal jerárquico}

Sea $g_{ij}$ el número de goles marcados por el equipo $i$ en el partido $j$.
El modelo base supone:
\begin{equation}
  g_{ij} \mid \lambda_{ij} \sim \text{Poisson}(\lambda_{ij}),
  \label{eq:poisson-marginal}
\end{equation}
con predictor log-lineal:
\begin{equation}
  \log(\lambda_{ij}) = \mu + \alpha_{k(i,j)} - \delta_{r(i,j)}
                        + \eta \cdot \mathbf{1}[\text{local}]_{ij}
                        + \beta \cdot \text{elo\_diff}_{ij},
  \label{eq:log-linear}
\end{equation}
donde $\mu$ es la intercepción global, $\alpha_k$ el efecto aleatorio de ataque
del equipo $k$, $\delta_k$ el efecto aleatorio de defensa, $\eta$ el coeficiente
de ventaja local y $\beta$ el coeficiente del diferencial de Elo normalizado.

Los efectos aleatorios por equipo se distribuyen jerárquicamente:
\begin{align}
  \alpha_k &\iid \mathcal{N}(0,\, \sigma_\alpha^2), \quad k = 1,\ldots,48,
  \label{eq:alpha-hier}\\
  \delta_k &\iid \mathcal{N}(0,\, \sigma_\delta^2), \quad k = 1,\ldots,48.
  \label{eq:delta-hier}
\end{align}
El supuesto de intercambiabilidad permite que equipos con pocas observaciones
directas tomen información de la distribución global, evitando sobreajuste
(regularización bayesiana jerárquica o \textit{partial pooling}).

\subsection{La corrección de Dixon-Coles}
\label{subsec:dc}

El modelo de Poisson con independencia marginal —ecuación~\eqref{eq:poisson-marginal}
con $g_{1j} \perp g_{2j} \mid \lambda_{1j}, \lambda_{2j}$— sobreestima
sistemáticamente las probabilidades de victorias por un gol (1-0, 0-1) y
subestima los empates de baja puntuación (0-0, 1-1), como documentan
\citet{dixon1997modelling} y \citet{karlis2003analysis}.
Este fenómeno refleja la dependencia táctica observada en fútbol:
los equipos ajustan su estrategia según el marcador vigente, reduciendo
la varianza conjunta por debajo de la que predicen dos Poisson independientes.

Para corregir este defecto manteniendo la tractabilidad, \citet{dixon1997modelling}
modifican la distribución conjunta del marcador mediante un factor de escala:
\begin{equation}
  P(g_1, g_2) = \text{Pois}(g_1;\, \lambda_1) \cdot \text{Pois}(g_2;\, \lambda_2)
               \cdot \tau(g_1, g_2,\, \lambda_1, \lambda_2,\, \rho),
  \label{eq:dc-joint}
\end{equation}
con función de corrección:
\begin{equation}
  \tau(g_1, g_2, \lambda_1, \lambda_2, \rho) = \begin{cases}
    1 - \lambda_1 \lambda_2 \rho & (g_1, g_2) = (0,0), \\
    1 + \lambda_2 \rho           & (g_1, g_2) = (1,0), \\
    1 + \lambda_1 \rho           & (g_1, g_2) = (0,1), \\
    1 - \rho                     & (g_1, g_2) = (1,1), \\
    1                            & g_1 + g_2 \geq 3.
  \end{cases}
  \label{eq:tau}
\end{equation}
El parámetro $\rho$ controla la dirección e intensidad de la corrección.
Valores negativos incrementan $P(0,0)$ y $P(1,1)$ y reducen $P(1,0)$ y $P(0,1)$.
Puede verificarse que $\sum_{g_1,g_2} P(g_1, g_2) = 1$ para todo $\rho$ en el
dominio admisible, garantizando una distribución de probabilidad válida.

\begin{proposition}
La log-verosimilitud ponderada del modelo Dixon-Coles es:
\begin{equation}
  \ell(\boldsymbol{\theta}) = \sum_{j=1}^{N} p_j \Bigl[
    \log \operatorname{Pois}(g_{1j};\, \lambda_{1j}) +
    \log \operatorname{Pois}(g_{2j};\, \lambda_{2j}) +
    \log \tau(g_{1j}, g_{2j}, \lambda_{1j}, \lambda_{2j}, \rho)
  \Bigr],
  \label{eq:loglik}
\end{equation}
donde $p_j$ son los pesos definidos en~\eqref{eq:temporal-weight}--\eqref{eq:tournament-mult}
y $\boldsymbol{\theta} = (\mu, \beta, \eta, \rho, \{\alpha_k\}, \{\delta_k\}, \sigma_\alpha, \sigma_\delta)$.
\end{proposition}

\subsection{Especificación de priors}

Los priors son débilmente informativos, concentrando la masa en rangos
plausibles sin dominar la verosimilitud de los 1{,}935 partidos de entrenamiento:
\begin{align}
  \mu          &\sim \mathcal{N}(0{,}3;\; 0{,}5),
    \quad e^{0{,}3} \approx 1{,}35 \text{ goles por partido},
    \label{eq:prior-mu}\\
  \eta         &\sim \mathcal{N}(0{,}2;\; 0{,}3),
    \quad e^{0{,}2} \approx 1{,}22 \text{ (ventaja local)},
    \label{eq:prior-eta}\\
  \beta        &\sim \mathcal{N}(0;\; 0{,}3),
    \label{eq:prior-beta}\\
  \rho         &\sim \mathcal{N}(0;\; 0{,}1) \cdot \mathbf{1}_{[-0{,}5,\; 0{,}3]},
    \label{eq:prior-rho}\\
  \sigma_\alpha &\sim \operatorname{Exponential}(5),
    \label{eq:prior-sigma-a}\\
  \sigma_\delta &\sim \operatorname{Exponential}(5).
    \label{eq:prior-sigma-d}
\end{align}
Los priors exponenciales sobre las escalas jerárquicas son preferibles a
priors half-normal en este contexto: al no tener masa en cero evitan el
embudo de Neal (\textit{Neal's funnel}) que surge cuando los efectos aleatorios
son near-zero y el muestreador HMC debe explorar una geometría altamente
curvilínea \citep{gelman2013bayesian}.
El prior sobre $\rho$ se trunca en $[-0{,}5;\, 0{,}3]$ para garantizar
$\tau(\cdot) > 0$ con las tasas $\lambda$ típicas del fútbol internacional.

\subsection{Identificabilidad y parametrización no centrada}

La separación ataque-defensa está identificada porque los efectos $\alpha_k$
tienen media cero por construcción: la intercepción $\mu$ absorbe el nivel
medio de goles y el contraste $\alpha_{k} - \delta_{r}$ captura toda la
variación entre equipos.

Para evitar el embudo de Neal en el muestreo MCMC, se emplea la
\textit{parametrización no centrada} \citep{papaspiliopoulos2007general}:
\begin{equation}
  \alpha_k = \sigma_\alpha \cdot \tilde{\alpha}_k,
  \qquad \tilde{\alpha}_k \iid \mathcal{N}(0, 1),
  \label{eq:noncentered}
\end{equation}
y análogamente para $\delta_k$.
Esta reparametrización desacopla la información sobre la escala ($\sigma_\alpha$)
de la información sobre los efectos individuales ($\tilde{\alpha}_k$),
produciendo geometrías del posterior más regulares y mejores tasas de aceptación.

%% ==========================================================
\section{Inferencia Bayesiana}
\label{sec:inference}
%% ==========================================================

\subsection{Implementación en Stan}

El modelo se implementa en el lenguaje probabilístico Stan
\citep{carpenter2017stan}.
El factor de corrección de Dixon-Coles se incorpora como término adicional
en el bloque \texttt{model} mediante la instrucción \texttt{target +=}:

\begin{verbatim}
for (n in 1:N) {
  real lam1 = exp(alpha + atk[team1[n]] - def[team2[n]] + ...);
  real lam2 = exp(alpha + atk[team2[n]] - def[team1[n]] + ...);
  real log_tau;
  if      (g1[n]==0 && g2[n]==0) log_tau = log(1 - lam1*lam2*rho);
  else if (g1[n]==1 && g2[n]==0) log_tau = log(1 + lam2*rho);
  else if (g1[n]==0 && g2[n]==1) log_tau = log(1 + lam1*rho);
  else if (g1[n]==1 && g2[n]==1) log_tau = log1m(rho);
  else                            log_tau = 0.0;
  target += peso[n] * (poisson_lpmf(g1[n]|lam1)
                     + poisson_lpmf(g2[n]|lam2)
                     + log_tau);
}
\end{verbatim}

Los pesos $p_j$ se aplican multiplicando directamente el aporte de cada
observación al \texttt{target}, equivalente a la verosimilitud ponderada
de la ecuación~\eqref{eq:loglik}.

\subsection{Diagnósticos de convergencia}

Se corren \textbf{4 cadenas} independientes con \textbf{3{,}000 iteraciones}
cada una (1{,}000 de calentamiento), resultando en \textbf{8{,}000 muestras
efectivas totales}. Se emplean \texttt{adapt\_delta = 0{,}99} y
\texttt{max\_treedepth = 12} para mejorar la exploración en regiones
de alta curvatura del posterior.

\begin{table}[H]
\centering
\caption{Diagnósticos MCMC para los parámetros del modelo.}
\label{tab:diagnostics}
\begin{tabular}{lrrrr}
\toprule
\textbf{Parámetro} & $\hat{R}$ & $n_{\text{eff}}$ & \textbf{Media} & \textbf{SD} \\
\midrule
$\mu$ (intercepción)  & 1{,}001 & 9{,}654  &  0{,}141 & 0{,}062 \\
$\beta$ (Elo)         & 1{,}001 & 11{,}328 &  0{,}510 & 0{,}097 \\
$\eta$ (local)        & 1{,}001 &  9{,}992 &  0{,}165 & 0{,}093 \\
$\rho$ (Dixon-Coles)  & 1{,}001 & 12{,}226 & $-$0{,}071 & 0{,}070 \\
$\sigma_\alpha$       & 1{,}002 &  5{,}091 &  0{,}062 & 0{,}049 \\
$\sigma_\delta$       & 1{,}001 &  4{,}036 &  0{,}069 & 0{,}053 \\
\midrule
\multicolumn{5}{l}{\textit{Divergencias: 2 de 8{,}000 iteraciones (0{,}025\%).}} \\
\bottomrule
\end{tabular}
\end{table}

El estadístico $\hat{R}$ de \citet{gelman1992inference} permanece por debajo
de 1{,}002 en todos los parámetros, muy por debajo del umbral de preocupación
habitual de 1{,}01 \citep{vehtari2021rank}.
Los valores de $n_{\text{eff}}$ superan los 4{,}000 en todos los casos.
Las 2 divergencias observadas (0{,}025\% del total) son estadísticamente
negligibles.

\subsection{Resumen del posterior}

\begin{table}[H]
\centering
\caption{Resumen del posterior de parámetros fijos (IC al 95\%).}
\label{tab:posterior}
\rowcolors{2}{lightgray}{white}
\begin{tabular}{lrrrrr}
\toprule
\textbf{Parámetro} & \textbf{Media} & \textbf{SD} & \textbf{2{,}5\%} & \textbf{50\%} & \textbf{97{,}5\%} \\
\midrule
$\mu$              &  0{,}141 & 0{,}062 &  0{,}018 &  0{,}141 &  0{,}262 \\
$\beta$ (Elo)      &  0{,}510 & 0{,}097 &  0{,}316 &  0{,}511 &  0{,}696 \\
$\eta$ (local)     &  0{,}165 & 0{,}093 & $-$0{,}021 &  0{,}165 &  0{,}345 \\
$\rho$ (DC)        & $-$0{,}071 & 0{,}070 & $-$0{,}206 & $-$0{,}072 &  0{,}065 \\
$\sigma_\alpha$    &  0{,}062 & 0{,}049 &  0{,}003 &  0{,}054 &  0{,}178 \\
$\sigma_\delta$    &  0{,}069 & 0{,}053 &  0{,}004 &  0{,}059 &  0{,}196 \\
\bottomrule
\end{tabular}
\end{table}

\paragraph{Efecto Elo ($\hat{\beta} = 0{,}510$).}
Un diferencial de Elo de 400 puntos multiplica la tasa esperada de goles del
equipo favorito por $e^{0{,}510} \approx 1{,}67$.
El IC al 95\% excluye el cero con holgura, confirmando el poder predictivo
del historial de rendimiento histórico.

\paragraph{Ventaja local ($\hat{\eta} = 0{,}165$).}
El IC al 95\% incluye marginalmente el cero ($[-0{,}021;\, 0{,}345]$), consistente
con la naturaleza neutral del fútbol internacional moderno, donde los torneos
de alto nivel frecuentemente se celebran en sedes neutrales.

\paragraph{Escalas jerárquicas ($\hat{\sigma}_\alpha \approx 0{,}062$, $\hat{\sigma}_\delta \approx 0{,}069$).}
Los valores reducidos reflejan que, una vez controlados el nivel medio $\mu$ y
el diferencial de Elo $\beta$, la heterogeneidad residual entre equipos en
ataque y defensa es modesta, plausible dada la homogeneización de la muestra
hacia equipos clasificados al Mundial.

%% ==========================================================
\section{Comparación de Modelos}
\label{sec:model-comparison}
%% ==========================================================

Durante el desarrollo del pipeline se evaluaron tres especificaciones:

\begin{enumerate}
  \item \textbf{Poisson independiente (sin DC).}
        La omisión del factor $\tau$ sobreestima $P(1{,}0)$ y $P(0{,}1)$
        en aproximadamente 8--12\% y subestima $P(0{,}0)$ y $P(1{,}1)$
        en magnitudes similares.
        El efecto se propaga en la simulación: el modelo sobreestima
        victorias ajustadas y subestima empates en fase de grupos.

  \item \textbf{Binomial Negativa independiente (sin DC).}
        Reemplazar Poisson por Binomial Negativa mejora el ajuste a la
        sobredispersión marginal —especialmente en partidos de alta puntuación—
        pero no resuelve el problema de correlación en marcadores bajos que
        corrige $\tau$.
        Adicionalmente, el parámetro de dispersión $\phi$ adicional debilita
        la concentración modal de la distribución predictiva, reduciendo la
        precisión en la predicción del marcador más probable.

  \item \textbf{Poisson + Dixon-Coles (modelo propuesto).}
        La corrección reasigna masa de probabilidad direccionalmente —hacia empates
        y alejada de victorias por un gol— sin introducir parámetros de dispersión.
\end{enumerate}

\paragraph{Verificación predictiva posterior.}
Se generan 1{,}000 réplicas del conjunto de entrenamiento desde la distribución
predictiva posterior.
Las distribuciones marginales de goles $P(G = g)$ se comparan con frecuencias
empíricas: el modelo reproduce con exactitud la fracción de partidos 0-0
(6{,}8\% observado; modelo: $6{,}7 \pm 0{,}4\%$) y 1-1
(8{,}3\% observado; modelo: $8{,}1 \pm 0{,}5\%$).
El modelo sin corrección de Dixon-Coles subestima ambas frecuencias en
aproximadamente 1{,}2 puntos porcentuales.

%% ==========================================================
\section{Simulación del Torneo}
\label{sec:simulation}
%% ==========================================================

\subsection{Muestreo predictivo posterior}

Para cada escenario simulado $s = 1, \ldots, S$:
\begin{enumerate}
  \item Se extrae una muestra $\boldsymbol{\theta}^{(s)}$ del posterior MCMC.
  \item Se calculan $(\lambda_1^{(s)}, \lambda_2^{(s)})$ mediante~\eqref{eq:log-linear}
        con $\boldsymbol{\theta}^{(s)}$ (sede neutral: $\mathbf{1}[\text{local}] = 0$).
  \item Se propone $(g_1^*, g_2^*) \sim \operatorname{Pois}(\lambda_1^{(s)}) \otimes
        \operatorname{Pois}(\lambda_2^{(s)})$ y se aplica remuestreo por importancia
        con pesos $\tau(g_1^*, g_2^*, \ldots)$ para obtener el marcador corregido.
\end{enumerate}
Este procedimiento preserva la correlación posterior entre parámetros y resultados,
produciendo probabilidades de avance correctamente calibradas.

\subsection{Fase de grupos}

Para la fase de grupos se simulan $S = 100{,}000$ escenarios completos.
En cada escenario, los seis partidos de cada grupo se resuelven según el marcador
simulado. La clasificación sigue las reglas oficiales de la FIFA: puntos,
diferencia de goles, goles marcados, enfrentamiento directo y sorteo.
Los ocho mejores terceros clasifican según los criterios oficiales del Anexo C
del reglamento FIFA.

\subsection{Rondas eliminatorias}

Se simulan $S = 10{,}000$ escenarios de eliminatorias.
Los empates al término del tiempo reglamentario se resuelven mediante
penales modelados como $\text{Bernoulli}(0{,}5)$, reflejando la incertidumbre
máxima dado que los datos históricos de penales en torneos de élite no muestran
desequilibrios estadísticamente significativos entre equipos de nivel comparable.

\paragraph{Actualización en tiempo real.}
El pipeline descarga automáticamente el repositorio \texttt{martj42} y
actualiza \texttt{data/live/resultados\_reales.csv} con los resultados del
torneo. Los partidos ya disputados se tratan como fijos en la simulación;
solo los pendientes se simrolan, colapsando progresivamente la incertidumbre.

%% ==========================================================
\section{Resultados}
\label{sec:results}
%% ==========================================================

\subsection{Efectos por equipo y fase de grupos}

\begin{table}[H]
\centering
\caption{Efectos posteriores medios, rating Elo y probabilidad de título para los diez principales candidatos.}
\label{tab:teams}
\rowcolors{2}{lightgray}{white}
\begin{tabular}{lrrrr}
\toprule
\textbf{Equipo} & $\hat{\alpha}_k$ & $\hat{\delta}_k$ & \textbf{Elo} & $P(\text{título})$ \\
\midrule
España       &  0{,}089 & $-$0{,}087 & 2017 & 13{,}5\% \\
Argentina    &  0{,}076 & $-$0{,}071 & 1995 &  9{,}2\% \\
Francia      &  0{,}071 & $-$0{,}068 & 1985 &  8{,}1\% \\
Inglaterra   &  0{,}063 & $-$0{,}060 & 1970 &  7{,}4\% \\
Brasil       &  0{,}068 & $-$0{,}055 & 1962 &  7{,}1\% \\
Portugal     &  0{,}058 & $-$0{,}049 & 1948 &  5{,}3\% \\
Alemania     &  0{,}052 & $-$0{,}053 & 1943 &  4{,}8\% \\
Países Bajos &  0{,}049 & $-$0{,}047 & 1938 &  4{,}2\% \\
Marruecos    &  0{,}038 & $-$0{,}063 & 1918 &  3{,}7\% \\
Uruguay      &  0{,}041 & $-$0{,}044 & 1912 &  3{,}1\% \\
\bottomrule
\end{tabular}
\end{table}

\subsection{Probabilidades por ronda eliminatoria}

\begin{table}[H]
\centering
\caption{Probabilidades de alcanzar cada ronda (\%).}
\label{tab:knockout}
\rowcolors{2}{lightgray}{white}
\begin{tabular}{lrrrrr}
\toprule
\textbf{Equipo} & \textbf{Octavos} & \textbf{Cuartos} & \textbf{Semis} & \textbf{Final} & \textbf{Título} \\
\midrule
España       & 88{,}2 & 61{,}4 & 34{,}1 & 21{,}5 & 13{,}5 \\
Argentina    & 82{,}7 & 51{,}3 & 26{,}8 & 15{,}6 &  9{,}2 \\
Francia      & 80{,}9 & 47{,}8 & 24{,}5 & 14{,}3 &  8{,}1 \\
Inglaterra   & 79{,}1 & 45{,}2 & 22{,}8 & 13{,}2 &  7{,}4 \\
Brasil       & 77{,}8 & 43{,}6 & 21{,}3 & 12{,}4 &  7{,}1 \\
Portugal     & 74{,}3 & 38{,}2 & 17{,}9 &  9{,}8 &  5{,}3 \\
Alemania     & 73{,}1 & 36{,}4 & 16{,}5 &  8{,}9 &  4{,}8 \\
Países Bajos & 70{,}8 & 33{,}7 & 14{,}9 &  7{,}8 &  4{,}2 \\
\bottomrule
\end{tabular}
\end{table}

España emerge como favorita respaldada por los más altos efectos de ataque y
defensa del posterior y el mayor rating Elo entre los participantes.
El hecho de que los ocho candidatos de la Tabla~\ref{tab:knockout} sumen
solo $\approx 59{,}7\%$ de probabilidad de título ilustra la imprevisibilidad
intrínseca del formato de eliminación directa y la importancia de reportar
distribuciones completas en lugar de pronósticos puntuales.

%% ==========================================================
\section{Discusión}
\label{sec:discussion}
%% ==========================================================

\subsection{Interpretación de $\rho$ y la corrección de Dixon-Coles}

El valor estimado $\hat{\rho} = -0{,}071$ es consistente con la estimación
original de \citet{dixon1997modelling} para datos de la Premier League
($\hat{\rho} \approx -0{,}13$), aunque de menor magnitud, posiblemente
porque los partidos entre selecciones de alto nivel exhiben menor dependencia
táctica que los de ligas nacionales con mayor asimetría de fuerzas.

Para cuantificar el efecto en un partido equilibrado con
$\lambda_1 = \lambda_2 = 1{,}15$:
\begin{align}
  P_{\text{DC}}(0,0)  &\approx e^{-2{,}30} \cdot (1 + 1{,}15^2 \cdot 0{,}071) \approx 0{,}094, \\
  P_{\text{ind}}(0,0) &= e^{-2{,}30} \approx 0{,}085.
\end{align}
La corrección incrementa $P(0,0)$ en $\approx +1$ punto porcentual
(10{,}6\% relativo), con un efecto simétrico sobre $P(1,1)$.
Estos ajustes son modestos en términos absolutos pero cualitativamente
importantes en torneos donde la clasificación de grupo depende de diferencias
de un solo gol.

\subsection{El rol del Elo como covariable}

La estimación robusta $\hat{\beta} = 0{,}510$ establece que el historial
acumulado de rendimiento aporta poder predictivo incremental más allá de
los efectos aleatorios del período de entrenamiento.
La complementariedad es conceptualmente clara: los efectos $\alpha_k$ y
$\delta_k$ capturan el rendimiento reciente en el subconjunto WC-vs-WC
(2010--actualidad), mientras que el Elo integra toda la historia desde 1872,
incluyendo partidos contra equipos no clasificados al torneo que quedan fuera
del conjunto de entrenamiento.

\subsection{Limitaciones}

\begin{enumerate}
  \item \textbf{Penales equiprobables.}
        El supuesto $P(\text{victoria en penales}) = 0{,}5$ ignora
        posibles ventajas de equipos con mejores porteros o tiradores.
        Los datos históricos disponibles no permiten estimar estas
        diferencias con precisión estadística razonable.

  \item \textbf{Ausencia de datos de lesiones.}
        El modelo trata cada equipo como una entidad estática; la baja de
        un jugador clave no se refleja en tiempo real.

  \item \textbf{Independencia entre jornadas.}
        No se modelan rotaciones, fatiga ni motivación diferencial en
        partidos donde la clasificación ya está asegurada.

  \item \textbf{Estacionariedad de parámetros durante el torneo.}
        Los efectos $\alpha_k$ y $\delta_k$ se mantienen fijos en toda la
        simulación. Una extensión natural sería incorporar dinámica de estado
        latente en el espíritu de \citet{rue2000prediction}.

  \item \textbf{Heterogeneidad de sede.}
        La ventaja local $\eta$ es homogénea; no se modelan diferencias por
        altitud (Ciudad de México, $\sim$\,2{,}200\,m), clima ni distancia
        de viaje.
\end{enumerate}

%% ==========================================================
\section{Conclusiones}
\label{sec:conclusion}
%% ==========================================================

Se ha presentado un sistema de predicción bayesiana completo para el
FIFA World Cup 2026, basado en un modelo jerárquico de Poisson con corrección
de dependencia de Dixon-Coles, inferido mediante HMC-NUTS en Stan sobre
1{,}935 partidos internacionales ponderados temporalmente.

Las principales contribuciones son:
\begin{enumerate}
  \item La integración en un único marco bayesiano de efectos aleatorios
        jerárquicos de ataque y defensa, ratings Elo como covariable histórica,
        y corrección de dependencia de Dixon-Coles con parámetro $\rho$ estimado
        del posterior.
  \item La demostración de que la corrección de Dixon-Coles mejora
        cuantificablemente la calibración de frecuencias de resultados exactos
        frente al modelo de Poisson independiente y la Binomial Negativa,
        sin introducir parámetros que debiliten la concentración modal.
  \item La cuantificación robusta del efecto Elo
        ($\hat{\beta} = 0{,}510$, IC\,95\%: $[0{,}316;\, 0{,}696]$),
        que establece el historial de rendimiento como predictor genuino.
  \item Un sistema de actualización automática en tiempo real que recalibra
        las probabilidades de avance a medida que se conocen los resultados.
\end{enumerate}

Las predicciones sitúan a España, Argentina y Francia como los tres candidatos
más probables al título, con más del 40\% de la probabilidad restante distribuida
entre los 45 equipos restantes —subrayando la imprevisibilidad del formato de
eliminación directa y la importancia de reportar distribuciones completas en
lugar de pronósticos puntuales.

%% ==========================================================
\begin{thebibliography}{99}

\bibitem[Baio y Blangiardo(2010)]{baio2010bayesian}
Baio, G. y Blangiardo, M. (2010).
\newblock Bayesian hierarchical model for the prediction of football results.
\newblock \textit{Journal of Applied Statistics}, 37(2), 253--264.

\bibitem[Bürkner(2017)]{burkner2017brms}
Bürkner, P.-C. (2017).
\newblock \texttt{brms}: An {R} package for {B}ayesian multilevel models using {S}tan.
\newblock \textit{Journal of Statistical Software}, 80(1), 1--28.

\bibitem[Carpenter et al.(2017)]{carpenter2017stan}
Carpenter, B., Gelman, A., Hoffman, M.~D., Lee, D., Goodrich, B., Betancourt, M.,
Brubaker, M., Guo, J., Li, P. y Riddell, A. (2017).
\newblock {S}tan: A probabilistic programming language.
\newblock \textit{Journal of Statistical Software}, 76(1), 1--32.

\bibitem[Dixon y Coles(1997)]{dixon1997modelling}
Dixon, M.~J. y Coles, S.~G. (1997).
\newblock Modelling association football scores and inefficiencies in the
  football betting market.
\newblock \textit{Journal of the Royal Statistical Society: Series C (Applied
  Statistics)}, 46(2), 265--280.

\bibitem[Gelman y Rubin(1992)]{gelman1992inference}
Gelman, A. y Rubin, D.~B. (1992).
\newblock Inference from iterative simulation using multiple sequences.
\newblock \textit{Statistical Science}, 7(4), 457--472.

\bibitem[Gelman et al.(2013)]{gelman2013bayesian}
Gelman, A., Carlin, J.~B., Stern, H.~S., Dunson, D.~B., Vehtari, A. y Rubin,
  D.~B. (2013).
\newblock \textit{Bayesian Data Analysis} (3.ª~ed.).
\newblock CRC Press.

\bibitem[Groll et al.(2018)]{groll2018prediction}
Groll, A., Schauberger, G. y Tutz, G. (2018).
\newblock Prediction of major international soccer tournaments based on
  team-specific regularized {P}oisson regression.
\newblock \textit{Journal of Quantitative Analysis in Sports}, 11(2), 97--115.

\bibitem[Karlis y Ntzoufras(2003)]{karlis2003analysis}
Karlis, D. y Ntzoufras, I. (2003).
\newblock Analysis of sports data by using bivariate {P}oisson models.
\newblock \textit{Journal of the Royal Statistical Society: Series D (The
  Statistician)}, 52(3), 381--393.

\bibitem[Maher(1982)]{maher1982modelling}
Maher, M.~J. (1982).
\newblock Modelling association football scores.
\newblock \textit{Statistica Neerlandica}, 36(3), 109--118.

\bibitem[martj42(2024)]{martj42}
martj42 (2024).
\newblock \textit{International Football Results from 1872 to 2024}.
\newblock Repositorio GitHub: \url{https://github.com/martj42/international_results}.

\bibitem[Papaspiliopoulos et al.(2007)]{papaspiliopoulos2007general}
Papaspiliopoulos, O., Roberts, G.~O. y Sköld, M. (2007).
\newblock A general framework for the parametrization of hierarchical models.
\newblock \textit{Statistical Science}, 22(1), 59--73.

\bibitem[Rue y Salvesen(2000)]{rue2000prediction}
Rue, H. y Salvesen, Ø. (2000).
\newblock Prediction and retrospective analysis of soccer matches in a league.
\newblock \textit{Journal of the Royal Statistical Society: Series D (The
  Statistician)}, 49(3), 399--418.

\bibitem[Vehtari et al.(2021)]{vehtari2021rank}
Vehtari, A., Gelman, A., Simpson, D., Carpenter, B. y Bürkner, P.-C. (2021).
\newblock Rank-normalization, folding, and localization: An improved $\hat{R}$
  for assessing convergence of {MCMC}.
\newblock \textit{Bayesian Analysis}, 16(2), 667--718.

\end{thebibliography}

\end{document}
